\documentclass[11pt,a4paper]{article}
\usepackage[utf8]{inputenc}
\usepackage[T1]{fontenc}
\usepackage[french]{babel}
\usepackage{geometry}
\usepackage{hyperref}
\usepackage{booktabs}
\usepackage{enumitem}
\usepackage{xcolor}
\usepackage{titlesec}
\usepackage{fancyhdr}
\usepackage{graphicx}
\usepackage{amsmath}

\geometry{margin=2.5cm}
\hypersetup{colorlinks=true, linkcolor=blue, urlcolor=blue}

\pagestyle{fancy}
\fancyhf{}
\rhead{Actuarial Notebook Agent}
\lhead{Documentation Fonctionnelle}
\rfoot{Page \thepage}

\title{%
  \textbf{Actuarial Notebook Agent}\\[0.5em]
  \large Documentation Fonctionnelle\\[0.3em]
  \normalsize Version 2.0 — Mars 2026
}
\author{Direction Actuarielle}
\date{}

\begin{document}
\maketitle
\tableofcontents
\newpage

% ─────────────────────────────────────────────────────────────────────────────
\section{Présentation générale}
% ─────────────────────────────────────────────────────────────────────────────

\subsection{Objet}
L'\textbf{Actuarial Notebook Agent} est une application web permettant de construire
des \textbf{tables de mortalité d'expérience} à partir de données de portefeuille assuré.
Elle s'appuie sur une interface Streamlit à trois onglets et sur un agent LLM (modèle GPT)
qui exécute une méthodologie actuarielle validée, observe les résultats intermédiaires et
produit un rapport professionnel.

\subsection{Cas d'usage principaux}
\begin{enumerate}[leftmargin=2cm]
  \item \textbf{Exploration méthodologique} : exécuter le pipeline de référence sur des
        données synthétiques pour valider la méthode avant toute analyse réelle.
  \item \textbf{Analyse d'un portefeuille réel} : charger un fichier CSV et obtenir
        automatiquement la table d'expérience avec rapport.
  \item \textbf{Personnalisation du pipeline} : sélectionner et réordonner les notebooks
        pour adapter la méthodologie (e.g.\,ajouter une étape de correction d'expositions).
\end{enumerate}

\subsection{Périmètre fonctionnel}
\begin{itemize}
  \item Chargement de données portefeuille (CSV) ou génération synthétique (50 000 contrats).
  \item Contrôle qualité automatisé des données.
  \item Calcul des expositions et taux bruts de mortalité par âge entier.
  \item Lissage Whittaker-Henderson (paramètre $\lambda$ configurable).
  \item Visualisation graphique (taux bruts vs lissés, comparaison TH/TF 00-02).
  \item Calcul du SMR global et par tranche d'âge avec intervalle de confiance.
  \item Export CSV de la table d'expérience finale.
  \item Rapport narratif généré automatiquement par LLM.
\end{itemize}

% ─────────────────────────────────────────────────────────────────────────────
\section{Données en entrée}
% ─────────────────────────────────────────────────────────────────────────────

\subsection{Format du fichier portefeuille}
Le fichier CSV doit contenir une ligne par contrat avec les colonnes suivantes :

\begin{center}
\begin{tabular}{llll}
\toprule
\textbf{Colonne} & \textbf{Type} & \textbf{Obligatoire} & \textbf{Description} \\
\midrule
\texttt{id} & chaîne & Oui & Identifiant unique du contrat \\
\texttt{date\_naissance} & date & Oui & Date de naissance de l'assuré \\
\texttt{date\_entree} & date & Oui & Date d'entrée en observation \\
\texttt{date\_sortie} & date & Oui & Date de sortie (fin d'étude, décès, résiliation) \\
\texttt{cause\_sortie} & \texttt{'deces'|\!'autre'} & Oui & Cause de sortie \\
\texttt{sexe} & \texttt{'H'|\!'F'} & Oui & Sexe de l'assuré \\
\bottomrule
\end{tabular}
\end{center}

\subsection{Paramètres d'analyse}
Les paramètres suivants sont définis dans le notebook 01 et peuvent être modifiés
par l'agent lors de l'analyse d'un fichier réel :

\begin{center}
\begin{tabular}{lll}
\toprule
\textbf{Paramètre} & \textbf{Valeur par défaut} & \textbf{Description} \\
\midrule
\texttt{SEXE} & \texttt{'H'} & Sexe analysé (H=hommes, F=femmes) \\
\texttt{DATE\_FIN\_OBSERVATION} & 2023-12-31 & Fin de la période d'observation \\
\texttt{LAMBDA\_WH} & 100 & Paramètre de lissage Whittaker-Henderson \\
\bottomrule
\end{tabular}
\end{center}

% ─────────────────────────────────────────────────────────────────────────────
\section{Données en sortie}
% ─────────────────────────────────────────────────────────────────────────────

\subsection{Table de mortalité exportée}
Un fichier CSV est exporté sous la forme \texttt{table\_mortalite\_<portefeuille>\_<sexe>.csv}
avec les colonnes :

\begin{center}
\begin{tabular}{ll}
\toprule
\textbf{Colonne} & \textbf{Description} \\
\midrule
\texttt{age} & Âge entier (20 à 90) \\
\texttt{E\_x} & Exposition centrale (années-personnes) \\
\texttt{D\_x} & Décès observés \\
\texttt{q\_x\_brut} & Taux brut annuel ($1-e^{-D_x/E_x}$) \\
\texttt{q\_x\_lisse} & Taux lissé Whittaker-Henderson \\
\texttt{q\_x\_ref} & Taux de référence TH/TF 00-02 \\
\texttt{SMR\_local} & Ratio observé/attendu par âge \\
\bottomrule
\end{tabular}
\end{center}

\subsection{Rapport d'analyse}
L'interface affiche, pour chaque étape, un rapport narratif en Markdown formaté
par le LLM (sans affichage du code Python). Un résumé exécutif clôture l'analyse.

% ─────────────────────────────────────────────────────────────────────────────
\section{Méthodologie actuarielle}
% ─────────────────────────────────────────────────────────────────────────────

\subsection{Calcul des expositions}
La méthode des \textbf{années-personnes} (Poisson) est utilisée. Pour chaque assuré $i$
et chaque âge entier $x$, l'exposition est :

\[
  e_{i,x} = \min(\text{date\_sortie}_i,\, x+1) - \max(\text{date\_entree}_i,\, x) \quad \text{(en années)}
\]

L'exposition centrale agrégée est $E_x = \sum_i e_{i,x}$.

\subsection{Taux bruts}
Le taux central brut est $\hat{\mu}_x = D_x / E_x$, et le taux annuel approché est :
\[
  \hat{q}_x^{\text{brut}} = 1 - e^{-\hat{\mu}_x}
\]

\subsection{Lissage Whittaker-Henderson}
Le lisseur minimise (sur l'échelle log) :
\[
  Q = \sum_x E_x (\log \hat{q}_x - y_x)^2 + \lambda \sum_x (\Delta^2 y_x)^2
\]
Solution matricielle : $\hat{y} = (W + \lambda D^\top D)^{-1} W y$ où $D$ est la matrice
de différences secondes et $W = \text{diag}(E_x)$.

\subsection{SMR}
Le SMR (Standardised Mortality Ratio) mesure la surmortalité relative :
\[
  \text{SMR} = \frac{\sum_x D_x}{\sum_x E_x \cdot q_x^{\text{réf}}}
\]
L'intervalle de confiance à 95\,\% (approche Poisson) est :
$\text{SMR} \pm 1{,}96\sqrt{D_{\text{obs}}} / D_{\text{att}}$.

\subsection{Table de référence}
La table TH 00-02 (hommes) / TF 00-02 (femmes) est approximée par interpolation
log-linéaire entre les valeurs de référence publiées par l'Institut des Actuaires.

% ─────────────────────────────────────────────────────────────────────────────
\section{Flux de traitement}
% ─────────────────────────────────────────────────────────────────────────────

\subsection{Pipeline complet}
\begin{enumerate}
  \item \textbf{NB01} — Chargement des données (CSV réel ou synthétique).
  \item \textbf{NB02} — Contrôle qualité : suppression des lignes incohérentes.
  \item \textbf{NB03} — Calcul des expositions $E_x$ et décès $D_x$ par âge entier.
  \item \textbf{NB04} — Lissage Whittaker-Henderson des taux bruts.
  \item \textbf{NB05} — Visualisation graphique (taux bruts, lissés, référence).
  \item \textbf{NB06} — Calcul du SMR global et par décennie.
  \item \textbf{NB07} — Export CSV de la table finale.
\end{enumerate}

\subsection{Comportement de l'agent}
L'agent suit une architecture \textbf{ReAct} (Reason + Act) :
après chaque appel à l'outil \texttt{execute\_python}, il analyse la sortie,
détecte les erreurs, adapte le code si nécessaire, puis passe à l'étape suivante.
Un LLM séparé (gpt-4o-mini) formate chaque sortie brute en texte lisible.

% ─────────────────────────────────────────────────────────────────────────────
\section{Limitations et hypothèses}
% ─────────────────────────────────────────────────────────────────────────────

\begin{itemize}
  \item Le sexe est uniforme sur l'analyse (pas de table mixte en une passe).
  \item Les âges hors plage [20, 90] sont exclus de la table de lissage.
  \item La table de référence TH/TF 00-02 est une approximation paramétrique ;
        remplacer par la table officielle pour les calculs réglementaires.
  \item Le lissage WH est appliqué sur l'échelle log — les âges avec exposition
        nulle ou taux brut nul sont exclus et interpolés implicitement.
\end{itemize}

\end{document}